\documentclass{beamer}

\usetheme{simple}

\usepackage{lmodern}
\usepackage[scale=2]{ccicons}
\usepackage[utf8]{inputenc}
\usepackage[spanish]{babel}
\usepackage{minted}

\usemintedstyle{emacs}

% TODO:
% position adjustement
% change colours
% 

% Watermark background (simple theme)

\setwatermark{\includegraphics[height=8cm]{img/Haskell-Logo.png}}

\title{Introducción a Haskell}
\subtitle{}
\date{17 de mayo de 2017}
\author{Anabel Ovide\\Ignacio Ballesteros\\Luis Eduardo Bueso\\Santiago Cervantes}
\institute{\url{https://github.com/edububa/haskell_course}}

\titlegraphic{
  \includegraphics[width=2cm]{img/fiwhiteblack.png}\hspace*{4cm}~%
  \includegraphics[width=2cm]{img/logo_acm_cromo.png}% \hspace{2cm}~%
  % \includegraphics[width=2cm]{img/logo_acm_cromo.png}
}

\begin{document}

\maketitle

\frame{\frametitle{Índice}\tableofcontents}

\definecolor{bg}{rgb}{0.90,0.90,0.90} % code bg color

%%%%%%%%%%%%%%%%%%%%%%%%%%%%%%%%%%%%%%%%%%%%%%%%%%%%%%%%%%%%%%%%%%%%%%%%%%%%% 
% INTRODUCCION
%%%%%%%%%%%%%%%%%%%%%%%%%%%%%%%%%%%%%%%%%%%%%%%%%%%%%%%%%%%%%%%%%%%%%%%%%%%%% 
\section{Introducción}
\subsection{Instalación}
\begin{frame}[fragile]
  \frametitle{Introducción}
  \framesubtitle{Instalación}
  \begin{itemize}
  \item\texttt{Con gestores de paquetes}
    \begin{itemize}
    \item\texttt{Debian/Ubuntu}
      \begin{minted}[bgcolor=bg]{shell}
        sudo apt-get update
        sudo apt-get install ghc
      \end{minted}
    \item\texttt{Arch}
      \begin{minted}[bgcolor=bg]{shell}
        sudo pacman -S ghc
      \end{minted}
    \item\texttt{macOS}
      \begin{minted}[bgcolor=bg]{shell}
        brew install ghc
      \end{minted}
    \end{itemize}
  \item\texttt{Desde la web}\\
    \url{https://www.haskell.org/platform/mac.html}\\
    \url{https://www.haskell.org/platform/windows.html}\\
    \url{https://www.haskell.org/downloads/linux}
  \end{itemize}
\end{frame}

\subsection{Introducción Histórica}
\begin{frame}[fragile]          %TODO
  \frametitle{Introducción}
  \framesubtitle{Introducción Histórica}

\end{frame}

\subsection{¿Qué es Haskell?}
\begin{frame}{Intoducción}      %TODO
  \framesubtitle{¿Qué es Haskell?}

\end{frame}

%%%%%%%%%%%%%%%%%%%%%%%%%%%%%%%%%%%%%%%%%%%%%%%%%%%%%%%%%%%%%%%%%%%%%%%%%%%%% 
% PRIMEROS PASOS
%%%%%%%%%%%%%%%%%%%%%%%%%%%%%%%%%%%%%%%%%%%%%%%%%%%%%%%%%%%%%%%%%%%%%%%%%%%%% 
\section{Primeros Pasos}
\subsection{GHCi}
\begin{frame}[fragile]
  \frametitle{Primeros pasos}
  \framesubtitle{GHCi}
  \begin{itemize}
  \item\texttt{Así abrimos el intérprete de Haskell}
    \begin{minted}[bgcolor=bg]{shell}
      \$ ghci
      GHCi, version 8.0.2: http://www.haskell.org/ghc/
      :? for help
      Prelude>
    \end{minted}

  \item\texttt{Podemos escribir expresiones aritméticas y lógicas:}
    \begin{minted}[bgcolor=bg]{shell}
      Prelude> 2 + 2
      4
      Prelude> True && False
      False
    \end{minted}
  \end{itemize}
\end{frame}

\begin{frame}[fragile]
  \frametitle{Primeros pasos}
  \framesubtitle{GHCi}
  \begin{itemize}
  \item\texttt{Podemos llamar a funciones}
    \begin{minted}[bgcolor=bg]{bash}
      Prelude> 2 + 2
      4
      Prelude> True && False
      False
    \end{minted}

  \item\texttt{Errores}
    \begin{minted}[bgcolor=bg]{shell}
      Prelude> 2 + "hola"

      <interactive>:6:1: error:
      • No instance for (Num [Char]) arising from a use
      of ‘+’
      • In the expression: 2 + "hola"
      In an equation for ‘it’: it = 2 + "hola"
    \end{minted}
  \end{itemize}
\end{frame}

\begin{frame}[fragile]
  \frametitle{Primeros pasos}
  \framesubtitle{GHCi}
  \begin{itemize}
  \item\texttt{Expresiones útiles:}
    \begin{minted}[bgcolor=bg]{shell}
      Prelude> :t 5
      5 :: Num t => t
      Prelude> :t 2
      2 :: Num t => t
      Prelude> :t "hola"
      "hola" :: [Char]

      Prelude> :l introduccion.hs
      [1 of 1] Compiling Main
      ( introduccion.hs, interpreted )
      Ok, modules loaded: Main.
      *Main>
    \end{minted}
  \end{itemize}
\end{frame}

\subsection{Primeras funciones}
\begin{frame}[fragile]
  \frametitle{Primeros pasos}
  \framesubtitle{Primeras funciones}
  % \begin{block}{Función:}

  % \end{block}
  \texttt{Ejemplos:}
  \begin{minted}[bgcolor=bg]{haskell}
    doubleMe x = x + x

    doubleUs x y = doubleMe x + doubleMe y

    doubleSmallNumber x = (if x > 100 then x else x*2) + 1
  \end{minted}
\end{frame}

\subsection{Listas}
\begin{frame}[fragile]
  \frametitle{Primeros pasos}
  \framesubtitle{Listas}
  % \begin{block}{Listas:}
  % \end{block}
  \texttt{Creando listas en un fichero \verb~.hs~:}
  \begin{minted}[bgcolor=bg]{haskell}
    list = [1,2,3,4,5]
    list' = [1..5]
  \end{minted}
  \texttt{Ejemplos de listas en el GHCi}
  \begin{minted}[bgcolor=bg]{shell}
    Prelude> let list = [1,2,3,4,5]
    Prelude> let list1 = [1..5]
    Prelude> list == list1
    list == list1
    True
    Prelude>
  \end{minted}
\end{frame}

\begin{frame}[fragile]
  \frametitle{Primeros pasos}
  \framesubtitle{Listas}
  Las listas pueden concatenarse con el operador \verb~++~ y\\construirse con el operador \verb~:~
  \begin{minted}[bgcolor=bg]{shell}
    Prelude> [1..9] ++ [10..19]
    [1,2,3,4,5,6,7,8,9,10,11,12,13,14,15,16,17,18,19]
    Prelude> "hello " ++ "world!"
    "hello world!"
    Prelude> 1:[2..5]
    1:[2..5]
    [1,2,3,4,5]
    Prelude> 'h':"ello"
    "hello"
    Prelude> 1:[]
    [1]
  \end{minted}
\end{frame}

\begin{frame}[fragile]
  \frametitle{Primeros pasos}
  \framesubtitle{Listas}
  \texttt{Las operaciones básicas en las listas son:}
  \begin{itemize}
  \item \verb~head~:
    \begin{minted}[bgcolor=bg]{shell}
      Prelude> head [1..10]
      1
    \end{minted}
  \item \verb~tail~:
    \begin{minted}[bgcolor=bg]{shell}
      Prelude> tail [1..10]
      [2,3,4,5,6,7,8,9,10]
    \end{minted}
  \item \verb~last~:
    \begin{minted}[bgcolor=bg]{shell}
      Prelude> last [1..10]
      10
    \end{minted}

  \end{itemize}
\end{frame}

\begin{frame}[fragile]
  \frametitle{Primeros pasos}
  \framesubtitle{Listas}
  \begin{itemize}
  \item \verb~init~:
    \begin{minted}[bgcolor=bg]{shell}
      Prelude> init [1..10]
      [1,2,3,4,5,6,7,8,9]
    \end{minted}
  \item \verb~length~:
    \begin{minted}[bgcolor=bg]{shell}
      Prelude> length [1..10]
      10
    \end{minted}
  \item \verb~reverse~:
    \begin{minted}[bgcolor=bg]{shell}
      Prelude> reverse [1..10]
      [10,9,8,7,6,5,4,3,2,1]
    \end{minted}
  \item \verb~take~:
    \begin{minted}[bgcolor=bg]{shell}
      Prelude> take 4 [1..10]
      [1,2,3,4]
    \end{minted}
  \end{itemize}
\end{frame}

\subsection{Tuplas}
\begin{frame}[fragile]
  \frametitle{Primeros pasos}
  \framesubtitle{Tuplas}
  % \begin{block}{Tuplas:}
  % \end{block}
\end{frame}

%%%%%%%%%%%%%%%%%%%%%%%%%%%%%%%%%%%%%%%%%%%%%%%%%%%%%%%%%%%%%%%%%%%%%%%%%%%%% 
% TIPOS Y TYPECLASSES
%%%%%%%%%%%%%%%%%%%%%%%%%%%%%%%%%%%%%%%%%%%%%%%%%%%%%%%%%%%%%%%%%%%%%%%%%%%%% 
\section{Tipos y Typeclasses}
\begin{frame}{Tipos y \emph{Typeclasses}}      %TODO

\end{frame}

%%%%%%%%%%%%%%%%%%%%%%%%%%%%%%%%%%%%%%%%%%%%%%%%%%%%%%%%%%%%%%%%%%%%%%%%%%%%% 
% SINTAXIS EN FUNCIONES
%%%%%%%%%%%%%%%%%%%%%%%%%%%%%%%%%%%%%%%%%%%%%%%%%%%%%%%%%%%%%%%%%%%%%%%%%%%%% 
\section{Sintaxis en funciones}
\subsection{Pattern matching}
\begin{frame}[fragile]
  \frametitle{Sintaxis en funciones}
  \framesubtitle{\emph{Pattern matching}}
  \begin{minted}[bgcolor=bg]{haskell}
    sayMe :: (Integral a) => a -> String
    sayMe 1 = "One!"
    sayMe 2 = "Two!"
    sayMe 3 = "Three!"
    sayMe 4 = "Four!"
    sayMe 5 = "Five!"
    sayMe x = "Not between 1 and 5"
  \end{minted}

  \begin{minted}[bgcolor=bg]{shell}
    Prelude> :l introduccion.hs
    [1 of 1] Compiling Main
    ( introduccion.hs, interpreted )
    Ok, modules loaded: Main.
    *Main> sayMe 5
    "Five!"
  \end{minted}
\end{frame}

\subsection{Guards}
\begin{frame}[fragile]
  \frametitle{Sintaxis en funciones}
  \framesubtitle{\emph{Guards}}
  \begin{minted}[bgcolor=bg]{haskell}
    compare :: (Ord a) => a -> a -> Ordering
    a `compare` b
       | a > b     = GT
       | a == b    = EQ
       | otherwise = LT
  \end{minted}
  \begin{minted}[bgcolor=bg]{shell}
    Prelude> :l introduccion.hs
    [1 of 1] Compiling Main
    ( introduccion.hs, interpreted )
    Ok, modules loaded: Main.
    *Main> 1 `myCompare` 2
    LT
    *Main> 1 `myCompare` 0.5
    GT
  \end{minted}
\end{frame}

\subsection{where}
\begin{frame}[fragile]
  \frametitle{Sintaxis en funciones}
  \framesubtitle{\emph{where}}
  \begin{minted}[bgcolor=bg]{haskell}
    initials :: String -> String -> String  
    initials firstname lastname = [f] ++ ". " ++ [l] ++ "."  
        where (f:_) = firstname  
              (l:_) = lastname    
  \end{minted}
\end{frame}
\subsection{let}
\begin{frame}[fragile]
  \frametitle{Sintaxis en funciones}
  \framesubtitle{\emph{let}}
  \begin{minted}[bgcolor=bg]{haskell}
    cylinder :: (RealFloat a) => a -> a -> a  
    cylinder r h = 
        let sideArea = 2 * pi * r * h  
            topArea = pi * r ^2  
        in  sideArea + 2 * topArea  
  \end{minted}
\end{frame}
\subsection{Case Expressions}
\begin{frame}[fragile]
  \frametitle{Sintaxis en funciones}
  \framesubtitle{\emph{Case Expressions}}
  \begin{minted}[bgcolor=bg]{haskell}
    describeList :: [a] -> String  
    describeList xs = "The list is " ++
      case xs of [] -> "empty."  
                 [x] -> "a singleton list."   
                 xs -> "a longer list."  
  \end{minted}
\end{frame}
%%%%%%%%%%%%%%%%%%%%%%%%%%%%%%%%%%%%%%%%%%%%%%%%%%%%%%%%%%%%%%%%%%%%%%%%%%%%% 
% RECURSION
%%%%%%%%%%%%%%%%%%%%%%%%%%%%%%%%%%%%%%%%%%%%%%%%%%%%%%%%%%%%%%%%%%%%%%%%%%%%% 
\section{Recursión}
\begin{frame}{Recursión}      %TODO

\end{frame}
\begin{frame}{Aclaraciones}
  Todo es mentira.
\end{frame}
\end{document}
